\src{170}
\begin{equation}\tag{1.3}
  \symup{H^{p}(X,B)=(R^{p}\Gamma)(B)=H^{p}\bigl(\Gamma(C(B))\bigr)}
\end{equation}
On a évidemment $\symup{h_A=\Gamma\hsh_A}$ ; d'autre part on vérifie que $\symup{\hsh_A}$
transforme objets injectifs en objets acycliques pour $\symup{\Gamma}$. On en conclut de
façon bien connue :

\result{PROPOSITION}{1} \uline{Il existe pour tout} $\symup{\Osh}$-\uline{Module} $\symup{A}$ \uline{un foncteur
spectral cohomologique sur} $\symup{\Ccat^{\Osh}}$, \uline{aboutissant au foncteur gradué}
$\symup{\bigl(\operatorname{Ext}_{\Osh}^{*}(X;A,B)\bigr)}$, \uline{et dont le terme initial est}
\begin{equation}\tag{1.4}
  \symup{E_{2}^{p,q}(A,B)=H^{p}\bigl(X,\Extsh_{\Osh}^{q}(A,B)\bigr)}
\end{equation}
On en déduit des \fquote{\uline{edge-homomorphisms}} et une suite exacte à cinq termes
que nous n'écrirons pas.

\result{COROLLAIRE}{1} \uline{Si} $\symup{A}$ \uline{est localement isomorphe à} $\symup{\Osh^{n}}$, \uline{alors on a
des isomorphismes canoniques}
\begin{equation}\tag{1.5}
  \symup{\operatorname{Ext}_{\Osh}^{p}(X;A,B)
  \overset{\sim}{\longleftarrow}
  H^{p}\bigl(X,\Homsh_{\Osh}(A,B)\bigr)}
\end{equation}
(donnés par des edge-homomorphisms de la suite spectrale). \uline{En particulier on a
un isomorphisme canonique}
\begin{equation}\tag{1.6}
  \symup{\operatorname{Ext}_{\Osh}^{p}(X;\Osh,B)=H^{p}(X,B)}
\end{equation}

Pour utiliser ces résultats, il faut savoir expliciter les
$\symup{\Extsh_{\Osh}^{p}(A,B)}$. Ce sont des foncteurs qui se calculent localement,
i.e. si $\symup{U}$ est un ouvert de $\symup{X}$, on a
\[
  \symup{\Extsh_{\Osh}^{p}(A,B)|U
  =\Extsh_{\Osh|U}^{p}(A|U,B|U)}
\]
comme il résulte du fait que la restriction à $\symup{U}$ d'un $\symup{\Osh}$-Module injectif
est un $\symup{(\Osh|U)}$-Module injectif. De plus, on a des homomorphismes fonctoriels,
pour $\symup{x\in X}$ fixé
\begin{equation}\tag{1.7}
  \symup{\Homsh_{\Osh}(A,B)_x
  \longrightarrow
  \operatorname{Hom}_{\Osh_x}(A_x,B_x)}
\end{equation}
qui se prolongent de façon unique en un homomorphisme de foncteurs cohomologiques
en $\symup{B}$ :
\begin{equation}\tag{1.8}
  \symup{\Extsh_{\Osh}^{p}(A,B)_x
  \longrightarrow
  \operatorname{Ext}_{\Osh_x}^{p}(A_x,B_x)}
\end{equation}

\result{PROPOSITION}{2} \uline{Si} $\symup{A}$ \uline{est, dans un voisinage de} $\symup{x}$, \uline{isomorphe au
conoyau d'un homomorphisme} $\symup{\Osh^{m}\longrightarrow\Osh^{n}}$, \uline{alors} (1.7) \uline{est un isomorphisme pour tout} $\symup{p}$. \uline{Il en est ainsi en
particulier si} $\symup{A}$ \uline{est un} $\symup{\Osh}$-\uline{Module cohérent} [3].

